
\needspace{2\baselineskip}
\color{uniblue}{\subsection{Логика отключения выключателя стороны НН трансформатора (ЛО НН)}}
\color{black}

\begin{enumerate}

%%%%%%%%%%%%%%%%%%%%%%%%%%%%%%%%%%%%%%%%%%%%%%%%%%%%%%%%%%%%%%%%%%%%% ОБЩИЕ СВЕДЕНИЯ %%%%%%%%%%%%%%%%%%%%%%%%%%%%%%%%%%%%%%%%%%%%%%%%%%%%%%%%%%%%%%

\item Общие сведения

\begin{enumerate}

\item
Логика отключения выключателя стороны НН трансформатора служит для формирования команды отключения в схему выключателя ввода секций НН трансформатора и команды запрета АПВ в схему АПВ выключателя ввода секции НН от защит в составе устройства и внешних защит, которые действуют через устройство.

\item
В устройстве логика выключателя стороны НН трансформатора представлена функциональным блоком «ЛО НН». Функциональный блок «ЛО НН» состоит непосредственно из логики отключения («ЛО») и логики запрета АПВ («ЗАПВ»).

\item 
В составе устройства предусматриваются две логики отключения стороны НН трансформатора (для схем управления вводными выключателями секций НН1 и НН2): ЛО НН1 и ЛО НН2.

\item
ЛО НН может быть выведена из работы оперативно путем активации сигнала «ОВ~ЛО~НН».
\item
В таблице \ref{lonn:tbl1} приведены параметры, необходимые для настройки функций в составе ЛО НН.


\end{enumerate}

%%%%%%%%%%%%%%%%%%%%%%%%%%%%%%%%%%%%%%%%%%%%%%%%%%%%%%%%%%%%%%%% ЛОГИКА ОТКЛЮЧЕНИЯ ВЫКЛЮЧАТЕЛЯ НН %%%%%%%%%%%%%%%%%%%%%%%%%%%%%%%%%%%%%%%%%%%%%%%%%%

\item Логика отключения выключателя (ЛО)

\begin{enumerate}

\item
Логика отключения в составе функционального блока «ЛО НН» представлена функцией «ЛО».

\item
Ввод функции в работу на этапе параметрирования устройства осуществляется программным переключателем SGF1 «Ввод\_функции». Информация о введенной функции отображается выходным сигналом «ЛО НН / ЛО: Ввод».

\item
Функция может быть выведена из работы оперативно путем активации сигнала «ОВ ЛО НН / ЛО», что характеризуется активным состоянием сигнала «ЛО НН / ЛО: Оперативный вывод» на выходе функции.

\item
Функция принимает сигнал срабатывания от общей логики отключения трансформатора и формирует сигналы «ЛО НН / ЛО: Отключение» и «ЛО НН / ЛО: Отключить аварийно», сигнал «ЛО НН / ЛО: Отключить аварийно» является импульсным, время импульса определяется выдержкой времени таймера T1 (уставка «Тимп»).
\item
Алгоритм работы функции выполнен в соответствии с рисунком \ref{lonn:img1}. 

\vspace{2mm}
\begin{figure}[H]
\centering
\includegraphics[width=0.8\textwidth,height=0.8\textheight,keepaspectratio]{\fbpath/16. ЛО В НН Т 35 ДЗТ/_latex/img/1.pdf}
\captionof{figure}{Функциональная схема логики отключения выключателя}\label{lonn:img1}
\end{figure}


\end{enumerate}
\FloatBarrier % Все float'ы ДО этой строки будут размещены здесь.
%%%%%%%%%%%%%%%%%%%%%%%%%%%%%%%%%%%%%%%%%%%%%%%%%%%%%%%%%%%%%%%% ЛОГИКА ЗАПРЕТА АПВ НН %%%%%%%%%%%%%%%%%%%%%%%%%%%%%%%%%%%%%%%%%%%%%%%%%%

\item Логика запрета АПВ выключателя ввода стороны НН (ЗАПВ)

\begin{enumerate}

\item
Логика запрета АПВ в составе функционального блока «ЛО НН» представлена функцией «ЗАПВ».

\item
Ввод функции в работу на этапе параметрирования устройства осуществляется программным переключателем SGF1 «Ввод\_функции». Информация о введенной функции отображается выходным сигналом «ЛО НН / ЗАПВ: Ввод».

\item
Функция может быть выведена из работы оперативно путем активации сигнала «ОВ ЛО НН / ЗАПВ», что характеризуется активным состоянием сигнала «ЛО НН / ЗАПВ: Оперативный вывод» на выходе функции.

\vspace{2mm}
\begin{figure}[H]
\centering
\includegraphics[width=0.8\textwidth,height=0.8\textheight,keepaspectratio]{\fbpath/16. ЛО В НН Т 35 ДЗТ/_latex/img/2.pdf}
\captionof{figure}{Функциональная схема логики запрета АПВ}\label{lonn:img2}
\end{figure}


\item
Функция принимает сигнал запрета АПВ от общей логики запрета АПВ в составе общей логики отключения трансформатора ЛО Т и формирует сигнал «ЛО НН / ЗАПВ: Запрет АПВ».

\item
Алгоритм работы функции выполнен в соответствии с рисунком \ref{lonn:img2}.



\end{enumerate}

\end{enumerate}
\FloatBarrier % Все float'ы ДО этой строки будут размещены здесь.